% Part: computability
% Chapter: computability-theory
% Section: computable-sets

\documentclass[../../../include/open-logic-section]{subfiles}

\begin{document}

\olfileid{cmp}{thy}{cps}
\olsection{可计算集合}

我们可以将可计算性的概念从可计算函数推广到可计算集合：

\begin{defn}
  设 $S$ 为自然数集合。则 $S$ 是\emph{可计算的}，
  当且仅当其特征函数~$\Char{S}$ 可计算。换言之，
  $S$~可计算，当且仅当函数
\[
\Char{S}(x) =
\begin{cases}
1 & 若 $x \in S$ \\
0 & 否则
\end{cases}
\]
可计算。类似地，关系 $R(x_0, \dots, x_{k-1})$ 可计算，当且仅当其特征函数可计算。

可计算的集合和关系也称为\emph{可判定的}。
\end{defn}

\begin{explain}
注意，我们现在有多种可计算性概念：偏函数的、（全）函数的以及集合的。不要把它们混淆！{} \iftag{TMs}{计算偏函数的图灵机，在函数有定义的输入值上，会返回函数的输出值；而计算集合的图灵机，则会在判定输入值是否属于该集合后，返回 $1$ 或~$0$。}{}
\end{explain}

\end{document}